\documentclass[10pt]{article}
\usepackage[utf8]{inputenc}
\usepackage[T1]{fontenc}
\usepackage[spanish,es-noquoting,es-tabla]{babel}
\usepackage{lmodern}
\usepackage[a4paper,landscape,margin=1.6cm]{geometry}
\usepackage{longtable,array,booktabs}
\usepackage{enumitem}
\usepackage[hidelinks]{hyperref}

\setlength{\parindent}{0pt}
\setlength{\parskip}{4pt}
\newcolumntype{P}[1]{>{\raggedright\arraybackslash}p{#1}}
\setlist{nosep,leftmargin=*}

\begin{document}

\begin{center}
{\Large\bfseries Fase 0. Matriz de análisis de casos}\\[2pt]
Ciencia de Datos -- Unidad 2 (Learning Analytics) -- Trabajo aplicativo\\
Universidad de Cuenca
\end{center}

Para esta fase revisamos las cinco herramientas vistas en clase. Además de las diapositivas, nos apoyamos en los boletines del proyecto LALA (n.º 3, enero de 2020; n.º 4, agosto de 2020; n.º 5, abril de 2021) y en el manual de usuario de FlipMyLearning. En la última columna anotamos las limitaciones que nosotros identificamos al leer esos materiales, no solo las que se comentaron en clase.

{\footnotesize
\begin{longtable}{P{2.2cm} P{3.4cm} P{3.6cm} P{2.7cm} P{2.2cm} P{4.3cm} P{5.0cm}}
\toprule
\textbf{Herramienta} & \textbf{¿Qué decisión apoya y de quién?} & \textbf{¿Qué datos usa y de dónde salen?} & \textbf{Nivel de analítica} & \textbf{Tipo} & \textbf{¿Cómo la validaron? ¿Con qué evidencia?} & \textbf{Limitación o riesgo que identificamos}\\
\midrule
\endhead
\bottomrule
\endfoot

\textbf{OnTask}\par (U. de Edimburgo; piloto con la UFRPE, Brasil)
& Ayuda al \textbf{docente} a decidir qué retroalimentación personalizada envía a cada estudiante y en qué momento. El docente arma reglas que activan textos distintos según la situación de cada estudiante, lo que le permite dar retroalimentación individual en cursos grandes.
& Características académicas de los estudiantes, sus interacciones con las actividades del curso y su rendimiento (notas). Los datos los sube o conecta el propio docente desde el LMS y sus registros de evaluación.
& Descriptiva. La parte de \emph{qué hacer} la definen las reglas del docente, no un modelo predictivo.
& Enseñanza (apoya la práctica docente).
& Antes del piloto se aplicó una encuesta de percepción de la retroalimentación (23 ítems Likert de 7 puntos): respondieron 36 de 60 estudiantes (58\,\%). En el piloto, los docentes dedicaron entre 60 y 130 minutos por semana a escribir retroalimentación. Al final del semestre se encuestó a los estudiantes sobre impacto y utilidad ($n=48$, 42,8\,\% de respuesta): perciben que la retroalimentación fue oportuna y les ayudó a ajustar sus estrategias.
& \begin{itemize}
\item La evidencia es de percepción (encuestas), no de rendimiento medido; tampoco hubo grupo de comparación.
\item Los propios estudiantes desconfían de la automatización: fue el ítem peor valorado ($M=4{,}97$ sobre 7).
\item Preocupa perder el lado afectivo de la retroalimentación.
\item Demanda tiempo al docente (1 a 2 horas semanales) y las reglas hay que rehacerlas en cada curso.
\item Si una regla está mal planteada, la herramienta no avisa.
\end{itemize}\\
\midrule

\textbf{iCoRa}\par (ESPOL)
& Ayuda al \textbf{estudiante} a decidir qué materias tomar el próximo semestre. La idea es que llegue a su sesión de consejería de 15 minutos con opciones ya analizadas, para usar ese tiempo en discutirlas con el consejero.
& Historial académico del estudiante sobre su malla (materias aprobadas, reprobadas y repetidas, con sus notas) e información histórica de cada materia: créditos, carga semanal, estimadores de dificultad, distribución de calificaciones y rendimiento de cohortes anteriores. Todo sale del sistema académico de ESPOL.
& Predictiva y prescriptiva: predice un rango de nota para cada materia y un promedio esperado para el conjunto elegido, y explica las variables que influyen.
& Aprendizaje (el usuario es el estudiante); apoya también la consejería.
& Según el boletín 4, en 2020 la herramienta estaba en pruebas de usuario de usabilidad y utilidad; el boletín 5 (2021) indica que pasaría a una nueva etapa de pruebas. No reportan todavía resultados. Su antecedente directo, las nuevas visualizaciones de SiCa (Sistema de Consejerías Académicas), sí se evaluaron: pretest a 175 consejeros, 187 capacitados, SUS, postest a 128 y uso real con 315 consejeros y 7714 estudiantes en el primer semestre de 2019.
& \begin{itemize}
\item Aún no hay evidencia de que mejore las decisiones o el rendimiento.
\item El estudiante podría usar la predicción para elegir solo materias \emph{fáciles} y retrasar su avance.
\item Un rango de nota puede leerse como una certeza; las explicaciones con fórmulas pueden no entenderse.
\item Al basarse en cohortes pasadas, repite sus patrones (por ejemplo, cambios de docente no se reflejan).
\item Puede restarle peso a la conversación con el consejero.
\end{itemize}\\
\midrule

\textbf{Sistema de consejería UCuenca (AvAc)}\par (LALA, inspirado en LISSA de KU Leuven)
& Ayuda al \textbf{docente consejero} (y a directores de carrera) a orientar al estudiante sobre su trayectoria y las materias que conviene tomar. Las sesiones se hacen a mitad de ciclo y antes de la matrícula.
& Registros académicos del sistema de la universidad: materias aprobadas, reprobadas, repetidas, anuladas y retiradas; notas; estructura de la malla con prerrequisitos; carga por ciclo; comparación con compañeros. Luego se le integró un panel de predicción de abandono.
& Descriptiva (trayectoria y rendimiento); con el panel integrado, también predictiva.
& Gestión académica (consejería); apoya la enseñanza.
& Piloto en cinco fases (preparación, acuerdo, capacitación, uso y mejora). En la fase de acuerdo participaron 75 docentes y funcionarios de 11 carreras. En la fase de uso hubo 31 consejeros activos: revisaron 522 estudiantes y 178 asistieron a consejería por invitación (6 más por iniciativa propia). Solo 25 estudiantes respondieron la encuesta, en general con opiniones positivas (Guerra et al., 2020).
& \begin{itemize}
\item La UCuenca no tenía un proceso de consejería: se introdujo a la vez la herramienta y el proceso, con resistencia por la carga de trabajo y sin horas asignadas.
\item El piloto empezó con docentes \emph{entusiastas}, lo que sesga los resultados a favor.
\item Solo 25 de 184 estudiantes atendidos respondieron la encuesta.
\item No se midió el efecto en el rendimiento o la permanencia.
\item A los estudiantes les cuesta interpretar su propia situación en el tablero.
\end{itemize}\\
\midrule

\textbf{Flip\-My\-Lear\-ning}\par (Universidad de Cuenca; plugin de Moodle basado en NoteMyProgress)
& Ayuda al \textbf{docente} a seguir semana a semana el avance de sus estudiantes frente a lo planificado y a decidir a quién contactar (permite escribir al estudiante o a su grupo). Al \textbf{estudiante} le ayuda a autorregularse.
& Registros de Moodle: tiempo invertido, sesiones (interacciones separadas por menos de 30 minutos), días activos, accesos y finalización de recursos y actividades, entrega de tareas, calificaciones, intentos en cuestionarios e interacciones con otros estudiantes. El docente configura las semanas y las horas esperadas de trabajo.
& Descriptiva; la sección de deserción agrupa a los estudiantes por comportamiento (\emph{clustering}), un primer paso hacia lo predictivo.
& Aprendizaje (autorregulación) y enseñanza.
& Según su artículo de presentación (Aguilar, Maldonado-Mahauad y Sigua, 2021), docentes y estudiantes evaluaron las visualizaciones y las consideraron útiles. También se identificó que el tiempo invertido, las sesiones y los indicadores de profundidad cognitiva y amplitud social ayudan a reconocer grupos en riesgo de abandono. No se reporta el tamaño de la muestra.
& \begin{itemize}
\item Los grupos del \emph{clustering} no son una predicción validada: no hay etiqueta ni métrica de acierto.
\item Si el docente no configura la finalización de actividades, \emph{haber visto} un recurso cuenta como completarlo, y el progreso se sobrestima.
\item Más tiempo conectado no significa más aprendizaje, y la regla de 30 minutos para separar sesiones es arbitraria.
\item Exige configuración semanal al docente.
\item Muestra foto, nombre y correo del estudiante, sin indicar cómo se gestiona el consentimiento.
\end{itemize}\\
\midrule

\textbf{Sistema de predicción de abandono (UC3M)}\par (adaptado en ESPOL, UCuenca, UACh, UPS y PUC)
& En las carreras, ayuda a \textbf{consejeros y directores} a decidir a qué estudiantes asesorar y qué materias recomendarles. En los MOOC de la PUC, ayuda al \textbf{docente} a decidir a qué grupo de riesgo enviar un mensaje.
& En las carreras, registros académicos de cada institución (cada una usa los datos que tiene disponibles). En los MOOC, la actividad semanal y las tareas en Coursera.
& Predictiva: estima la probabilidad de abandono. En los MOOC agrupa a los estudiantes en tres niveles de riesgo y el modelo se recalcula cada semana.
& Gestión (carreras) y enseñanza (MOOC).
& Los pilotos midieron utilidad, impacto y rendimiento con los registros de uso y una encuesta voluntaria a directores y docentes. A partir de los pilotos se reubicaron los paneles para hacerlos más accesibles y, en ESPOL, se amplió el conjunto de datos y se agregó información explicativa. El boletín 5 resume que el uso y la efectividad fueron buenos, \emph{aunque hay dificultades para medir este indicador}.
& \begin{itemize}
\item La efectividad casi no se pudo medir: no hay comparación entre estudiantes con y sin intervención.
\item La encuesta voluntaria tiende a recoger a los usuarios más satisfechos.
\item En ESPOL los consejeros solo ven la predicción durante la consejería al inicio del semestre, cuando ya es tarde para actuar dentro del ciclo.
\item Las primeras versiones no explicaban la predicción, lo que dificulta confiar en ella.
\item Cada institución usa datos distintos, así que los resultados no son comparables.
\end{itemize}\\
\end{longtable}
}

\section*{Lo que tomamos de estos casos para nuestra propuesta}

Al comparar las cinco herramientas vimos un patrón: casi todas apoyan la consejería al \emph{inicio} del semestre (AvAc, iCoRa, la predicción en ESPOL), pero ninguna alerta a tiempo \emph{durante} el ciclo usando la actividad de eVirtual. FlipMyLearning sí usa Moodle, pero está pensada para el docente de cada curso y no para el tutor que acompaña a un grupo de estudiantes. Esa es la brecha que queremos cubrir.

\begin{enumerate}
\item \textbf{De OnTask} tomamos la idea de que cada alerta debe venir con una acción concreta y fácil de ejecutar (por ejemplo, un mensaje con plantilla).
\item \textbf{De AvAc y de la predicción UC3M} tomamos el uso de datos que la universidad ya tiene y la integración con el proceso de consejería existente; nuestra herramienta podría complementar a AvAc.
\item \textbf{De FlipMyLearning} tomamos los indicadores de actividad en Moodle (sesiones, días activos), pero cuidando sus limitaciones: no contamos horas conectadas y comparamos a cada estudiante con su propio curso.
\item \textbf{La limitación común} es la falta de evidencia de impacto con un grupo de comparación. Por eso nuestro plan de validación propone comparar paralelos con y sin herramienta y registrar cada intervención.
\end{enumerate}

\section*{Fuentes}
\begin{itemize}
\item Proyecto LALA. Boletín informativo n.º 3 (enero de 2020), n.º 4 (agosto de 2020) y n.º 5 (abril de 2021).
\item Manual de usuario de FlipMyLearning, versión 1.0 (2021).
\item Aguilar, B., Maldonado-Mahauad, J. y Sigua, E. (2021). \emph{FlipMyLearning: A Tool for Monitoring and Predicting Learner Behavior in Moodle}.
\item Guerra, J. et al. (2020). Adaptation and evaluation of a learning analytics dashboard to improve academic support at three Latin American universities. \emph{British Journal of Educational Technology}, 51(4), 973--1001.
\item Diapositivas de la Unidad 2.
\end{itemize}

\end{document}
